% \iffalse meta-comment
%<*internal>
\iffalse
%</internal>
%<*readme>
----------------------------------------------------------------
template --- short description
E-mail: yannislaz@gmail.com
Released under the LaTeX Project Public License v1.3c or later
See http://www.latex-project.org/lppl.txt
----------------------------------------------------------------
This file provides a template for defining a class.
%</readme>
%<*todo>
add tcolorbox support
%</todo>
%<*internal>
\fi
\def\nameofplainTeX{plain}
\ifx\fmtname\nameofplainTeX\else
  \expandafter\begingroup
\fi
%</internal>
%<*install>
\input docstrip.tex
\keepsilent
\askforoverwritefalse
\preamble
----------------------------------------------------------------
template --- short description.
E-mail: yannislaz@gmail.com
Released under the LaTeX Project Public License v1.3c or later
See http://www.latex-project.org/lppl.txt
----------------------------------------------------------------

\endpreamble
\postamble
 Copyright (C) 2013 by Dr. Yiannis Lazarides <yannislaz@gmail.com>
\endpostamble
\usedir{tex/latex/\jobname}
\generate{
  \file{\jobname.sty}{\from{\jobname.dtx}{package}}
 }
%</install>
%<install>\endbatchfile

%<*internal>
\usedir{source/latex/\jobname}
\generate{
  \file{\jobname.ins}{\from{\jobname.dtx}{install}}
}
\nopreamble\nopostamble
\usedir{doc/latex/demopkg}
\generate{
  \file{README.txt}{\from{\jobname.dtx}{readme}}
}
\generate{
  \file{test-01.tex}{\from{\jobname.dtx}{test-01}}
}
\generate{
  \file{TODO.tex}{\from{\jobname.dtx}{TODO}}
}
\ifx\fmtname\nameofplainTeX
  \expandafter\endbatchfile
\else
  \expandafter\endgroup
\fi
%</internal>
%<*driver>
\documentclass{ltxdoc}
\EnableCrossrefs
\CodelineIndex
\RecordChanges
\begin{document}
\DocInput{HS.dtx}
\end{document}
%</driver>
% \fi
%
% %%%%%%%%%%%%%%%%%%%%%%%%%%%%%%%%%%%%%%%%%%%%%%%%%%%%%%%%%%%%%%%%%%%%
%
% \CheckSum{301}
%
% \CharacterTable
%  {Upper-case    \A\B\C\D\E\F\G\H\I\J\K\L\M\N\O\P\Q\R\S\T\U\V\W\X\Y\Z
%   Lower-case    \a\b\c\d\e\f\g\h\i\j\k\l\m\n\o\p\q\r\s\t\u\v\w\x\y\z
%   Digits        \0\1\2\3\4\5\6\7\8\9
%   Exclamation   \!     Double quote  \"     Hash (number) \#
%   Dollar        \$     Percent       \%     Ampersand     \&
%   Acute accent  \'     Left paren    \(     Right paren   \)
%   Asterisk      \*     Plus          \+     Comma         \,
%   Minus         \-     Point         \.     Solidus       \/
%   Colon         \:     Semicolon     \;     Less than     \<
%   Equals        \=     Greater than  \>     Question mark \?
%   Commercial at \@     Left bracket  \[     Backslash     \\
%   Right bracket \]     Circumflex    \^     Underscore    \_
%   Grave accent  \`     Left brace    \{     Vertical bar  \|
%   Right brace   \}     Tilde         \~}
%
% \GetFileInfo{HS.dtx}
%
% \title{The \textsf{HS} package\thanks{This file
%         has version number \fileversion, last
%         revised \filedate.}}
% \author{Dr Yannis Lazarides\\y.lazarides@habtoorspecon.com}
% \maketitle
%
% ^^A  The following were copied verbatim from source2e.tex.
% \DoNotIndex{\def,\long,\edef,\xdef,\gdef,\let,\global}
% \DoNotIndex{\if,\ifnum,\ifdim,\ifcat,\ifmmode,\ifvmode,\ifhmode,%
%             \iftrue,\iffalse,\ifvoid,\ifx,\ifeof,\ifcase,\else,\or,\fi}
% \DoNotIndex{\box,\copy,\setbox,\unvbox,\unhbox,\hbox,%
%             \vbox,\vtop,\vcenter}
% \DoNotIndex{\@empty,\immediate,\write}
% \DoNotIndex{\egroup,\bgroup,\expandafter,\begingroup,\endgroup}
% \DoNotIndex{\divide,\advance,\multiply,\count,\dimen}
% \DoNotIndex{\relax,\space,\string}
% \DoNotIndex{\csname,\endcsname,\@spaces,\openin,\openout,%
%             \closein,\closeout}
% \DoNotIndex{\catcode,\endinput}
% \DoNotIndex{\jobname,\message,\read,\the,\m@ne,\noexpand}
% \DoNotIndex{\hsize,\vsize,\hskip,\vskip,\kern,\hfil,\hfill,\hss}
% \DoNotIndex{\m@ne,\z@,\z@skip,\@ne,\tw@,\p@}
% \DoNotIndex{\dp,\wd,\ht,\vss,\unskip}
%
% ^^A  The following are specific to HS.dtx.
% \DoNotIndex{\@currenvir,\@gobble,\@gobblefour,\@ifundefined,\@makeother}
% \DoNotIndex{\@undefined,\active,\chardef,\day,\do,\dospecials,\E,\end,\I}
% \DoNotIndex{\if@tempswa,\L,\LaTeX,\loop,\MessageBreak,\month}
% \DoNotIndex{\newenvironment,\number,\repeat,\reserved@b,\reserved@c}
% \DoNotIndex{\two@digits,\year,\*,\^}
%
% ^^A  Define some commands to help delineate my changes.
% \newcommand{\startfcchanges}{^^A
%   \centerline{^^A
%     \makebox[0pt]{^^A
%       \raisebox{-2\baselineskip}[0pt][0pt]{^^A
%       \makebox[0pt][l]{\rule{1pt}{2\baselineskip}}}^^A
%       \rule{1em}{1pt}^^A
%       \rule{\linewidth}{1pt}^^A
%       \rule{1em}{1pt}^^A
%       \raisebox{-2\baselineskip}[0pt][0pt]{^^A
%       \makebox[0pt][r]{\rule{1pt}{2\baselineskip}}}^^A
%     }^^A
%   }^^A
%  \noindent
% }
% \newcommand{\stopfcchanges}{^^A
%   \centerline{^^A
%     \makebox[0pt]{^^A
%       \raisebox{0pt}[0pt][0pt]{^^A
%       \makebox[0pt][l]{\rule{1pt}{2\baselineskip}}}^^A
%       \rule{1em}{1pt}^^A
%       \rule{\linewidth}{1pt}^^A
%       \rule{1em}{1pt}^^A
%       \raisebox{0pt}[0pt][0pt]{^^A
%       \makebox[0pt][r]{\rule{1pt}{2\baselineskip}}}^^A
%     }^^A
%   }^^A
% }
%
% \begin{abstract}
%
%  This package offers a number of macros for typesetting  Contractual letters
%  It arose out of a need to provide something simpler and fullfilled my needs. It is
%  easily modifiable via a local definitions file or better fork it for your own purposes.
%  It keeps commands to a minimum to avoid information overload. It offers a number
%  of settings, in order to make it easily customizable.
% \end{abstract}
% \section{Introduction}
%
% Official correspondence, consists of a few constant fields; the package provides macros for
% these. The |date| the letter was drafted, the filing reference of the sender, the Reference label
% the signature, any enclosures and to whom the letter was also distributed.
%
% \DescribeMacro{Date} The command |\Date|\oarg{optional argument} is used to automatically typeset
% today's date, or the argument provided. If the file is used again in a future date the date is cached as
% the first time the file was processed. This way copies of old letters can be reprocessed.
%
% \DescribeMacro{\RE}
%  The command |\RE|\oarg{optional argument}\marg{text} is used to typeset the reference
%  on top of a letter.
% \paragraph{Sample usage}
% |HS| works much like any package by including it in your preamble using:
% \begin{verbatim}
% \includepackage{HS}
% \end{verbatim}
% I avoided any options to keep it simple for people not very familiar with \LaTeX\  to use and to modify.
% \StopEventually{^^A
%   \PrintChanges
%   \PrintIndex
% }
%
% \section{Implementation}
%
% Most users can stop reading at this point.  The Implementation section
% contains the annotated source code for the \textsf{HS}
% package itself, which is useful only to people who want a detailed and
% precise explanation of how \textsf{HS} works.
%
%
%    \begin{macrocode}
%<*package>
%    \end{macrocode}
% We first load some packages that we use. The |lipsum| is provided so that you can experiment
% while developing your own customization.
%
%    \begin{macrocode}
\RequirePackage{lastpage}
\RequirePackage{graphicx}
\RequirePackage{lipsum}
\RequirePackage{url}
\RequirePackage{color}
\DeclareUrlCommand\email{\color{blue}\urlstyle{sf}}
%    \end{macrocode}
%
% The letter pages are numbered in a specific style, prevalent in Construction Projects as Page 1 of 2 etc.
% It is to put it mildly not very wise to number the first page, but this is the house style. Try and convince
% corporate types that it is unecessary.
%
%    \begin{macrocode}
\def\ps@plain{%
	\let\@oddhead\@empty
	\def\@oddfoot{\hfill \footnotesize \sffamily Page \textsf{\thepage} of \pageref{LastPage}}%
	\def\@evenfoot{\hfill\footnotesize \sffamily Page \textsf{thepage} of \pageref{LastPage} }}
	\pagestyle{plain}
%    \end{macrocode}
%
% \begin{macro}{\Date}
% The command |\Date|\oarg[date]\marg{date} takes one argument
% and typesets the date. It leaves the correct amount of space 
% below it.
%    \begin{macrocode}
 \gdef\letterdate{\today}
\newcommand\thedate[1]{%
     \@ifundefined{auxdate}{%
        \gdef\letterdate{\par\leavevmode#1\par}}%
       {\gdef\letterdate{\par\leavevmode\auxdate\par}}%
}
\newcommand{\Date}[1][]{%
      \if!#1!  \thedate{\today}%
          \else
             \thedate{#1}%
       \fi
     \vspace{1\baselineskip}%
    \letterdate
}
%    \end{macrocode}
% \end{macro}
%
% This will be extended to provide a cache version later to enable re-runs without worrying about the
% date changing.
%
% \begin{macro}{\OurRef}
% \begin{macro}{\refprefix}
% The user command \cs{OurRef} is used to provide the automatic reference to a letter. The user
% command \cs{refprefix} provides the prefix such as HS-142/yl/1002
% The Reference Number is stored the a counter |total|.
%    \begin{macrocode}
\newcounter{total}
%    \end{macrocode}
%    \begin{macrocode}
\newcommand\refprefix{HS-142/HLG/YL/mr/}
%    \end{macrocode}
%
% \subsection{Auto-incrementing references}
% We keep track of an autoincrement counter by using two files. The first one 
% |autoincrementtotals.lcl| keeps the totals. If you have five files letter-01.tex, letter-02.tex
% this counter will keep the latest reference number. Don't wipe it out. 
% The second file relates to the |\jobname| and is |jobname.lcl|. It caches the tracking
% number for a particular letter.
%
% First we check if the totalizer files exists. In case this is our first letter, it will
% create the file.
%
%    \begin{macrocode}
\IfFileExists{autoincrementtotals.lcl}{}{\newwrite\tempfile
	\immediate\openout\tempfile=autoincrementtotals.lcl
	\immediate\write\tempfile{\thetotal}%
	\immediate\closeout\tempfile}
%    \end{macrocode}
%
% The auxiliary macro |\autoinc| reads from the |\jobname.lcl| the
% cached value of the counter and stores it in the command |\OurRef|. This is  user 
% command.
%
%    \begin{macrocode}
\def\autoinc{%
     	\newread\inputstream
     	\immediate\openin\inputstream=\jobname.lcl
     	\immediate\read\inputstream to \auxcommand
        \immediate\read\inputstream to \auxdate
     	\immediate\closeout\inputstream
        
   	\newcommand\OurRef{%
           \leavevmode\refprefix
           \ifnum9>\auxcommand 100\auxcommand\else
               \ifnum99>\auxcommand 10\auxcommand\else
                  \ifnum999>\auxcommand \auxcommand
                  \fi
               \fi
          \fi
           
          
           \vspace{1\baselineskip}%
       }
       
}
%    \end{macrocode}
%
% First we check if a file with the |\jobname.lcl| exists holding the RefNumber
% If it does not exist increment the totalizer and write to the values to both files.
%    \begin{macrocode}
\IfFileExists{\jobname.lcl}{\autoinc}{%
    \newread\inputstream
    \immediate\openin\inputstream=autoincrementtotals.lcl
    \immediate\read\inputstream to \auxcommand
    \immediate\closeout\inputstream
    \setcounter{total}{\auxcommand}        
    \stepcounter{total}
%    \end{macrocode}
%  Having cached the reference number, we now upadte the |autoincrementtotals.lcl| file.
%    \begin{macrocode}
	\newwrite\tempfile
	\immediate\openout\tempfile=autoincrementtotals.lcl
	\immediate\write\tempfile{\thetotal}%
	\immediate\closeout\tempfile

	\gdef\OurRef{\leavevmode\refprefix\thetotal\vspace{1\baselineskip}}
	%% write to jobname.lcl
	\newwrite\tempfile
	\immediate\openout\tempfile=\jobname.lcl
	\immediate\write\tempfile{\thetotal}%
        \immediate\write\tempfile{\letterdate}%has error here Date not known yet
	\immediate\closeout\tempfile
}
%    \end{macrocode}
% \end{macro} 
% \end{macro}
%
% \subsection{Adressee}
%
% Here again, the way the addressee is handled depends on the house-style. Customization
% is king here otherwise you will find you need to do a lot of typing.
%
% \begin{macro}{\TO} The macro |\TO|\oarg{Recipient details} takes one optional argument and
%  defaults to the most typical recipient.
%    \begin{macrocode}
\newcommand{\TO}[1][Test]{%
    \leavevmode\par
    Mr. Ziad Hoedroge\\
    Project Director\\
    M/s Habtoor Leighton Group\\
    Al Habtoor Engineering Enterprises Co., LLC.\\
    P.O.Box 2248\\
    Doha\\
    Qatar\\
    \par
    \vspace{2\baselineskip}%
    Dear Sir
   \vspace{\baselineskip}%
}
%    \end{macrocode}
% \end{macro}
%
% Define your own macros, if you regularly send letters to a limited number of recipients.
% A factory command can be provided in a future release.
% 
% \begin{macro}{\TOMCD}
%    \begin{macrocode}
\newcommand{\TOMCD}[1][Test]{%
    \leavevmode
    Mr. Mohammad Azzam\\
    Mr. Marios Charalambous\\
    M/s HLS-DSI JV\\
    Dubai\\
    UAE\\
    \\
    \textbf{BY HAND}
    \par
    \vspace{2\baselineskip}
    Dear Sir
    \vspace{\baselineskip}
}
\parindent0pt
\parskip1em
%    \end{macrocode}
% \end{macro}
% \begin{macro}{\RE}
% Typeset the subject. The subject is normally set in and the 
% subject is typeset in
% bold and underlined.
%    \begin{macrocode}
\newlength\tempindent
\settowidth\tempindent{Sub:\hspace{1em}}
%
\newcommand{\RE}[2][]{%
  \parskip0pt\par\noindent
  \hangindent\tempindent
  \textbf{#1#2}%
  }
%    \end{macrocode}
% \end{macro}
%
% \subsection{Signatures}
%
%  We provide some general commands first. The author has the option to print a letter both using
%  a signature from an image or to allow space for the author to sign on the printed letter. When I
%  email letters rather than deliver them in hard copies that is what I use.
%  We first define booleans to handle the two cases. We also provide a switch for the author. The
%  default is to typeset from an image.
%
% \begin{macro}{\@signatureimage}
% \begin{macro}{\NoSignatureImage} The conditional |\@signatureimage| is used to provide
% a switch for deciding if we need to print the signature image or not. The command |NoSignatureImage|
% is an author command to set it to true or false.
%
%    \begin{macrocode}
\newif\if@signatureimage
         \@signatureimagetrue
\def\NoSignatureImage{\@signatureimagefalse} 
%    \end{macrocode}
% \end{macro}
% \end{macro}
%
% \begin{macro}{\YL}
% Typesets my signature and title
%    \begin{macrocode}

\def\SetSignatureImage#1{#1}
\SetSignatureImage{\gdef\signatureimage{dryl.jpg}}
\newcommand{\YL}{%
    \leavevmode\par
    Yours faithfully,\\
    For HLS-DSI JV\par
    \if@signatureimage\includegraphics[width=4cm]{\signatureimage}\\
        \else
        \vspace*{1.8cm}
    \fi
    Dr. Yiannis Lazarides\\
    Habtoor City Project, Dubai\\
    Project Director\\
   \email{y.lazarides@hlspecon.com}\\
    \vskip1.3\baselineskip
}
%    \end{macrocode}
%  \end{macro}
%
% \begin{macro}{\cc} The macro |\copies| typesets recipient names for whom letters were 
% distributed.
%    \begin{macrocode}
\newcommand{\cc}[1][]{%
    \if!#1! 
       \leavevmode\par
       D\&S HO\\
       HLS HO\\
      \else#1
   \fi
}
%    \end{macrocode}
% \end{macro}
% 
% \begin{macro}{\encl} The command |\encl| typesets our typical method of indicating enclosures.
% again the macro has an optional command.
%
%    \begin{macrocode}
\newcommand{\encl}{%
    \leavevmode\par
Encl. as above 
}
\long\def\signature{%
\YL
\encl}

\AtBeginDocument{%
    \sffamily
    \mbox{}
    %\vspace*{-2.5cm}%
    \includegraphics[width=\textwidth]{JVlogo.JPG}%include in head better
   
}
\newenvironment{mcdletter}[1][]{
\Date\par 
\OurRef\par
\TOMCD\par
\textbf{#1}\par

Please obtain quotations and process the necessary documentation, for the above materials as early
as possible and provide us with your recommendation and price comparisons in order for us to release the
order.
}{\signature
\par}



%    \end{macrocode}

% \end{macro}
%
%    \begin{macrocode}
%</package>
%    \end{macrocode}
%
% \Finale
%
